%======================================================================
\section{Introduction}
\label{sec:intro}

The problem treated here is easily stated.  A listed option market quotes,
for one underlying asset and one expiry, the prices of calls and puts at a
few dozen strikes.  A fitter must produce from these quotes a complete curve
--- an implied volatility at every strike --- that reproduces the quotes to
within their own uncertainty, extends beyond them sensibly, and can be
differentiated, integrated, and compared across expiries by every system
downstream.  The problem looks like curve fitting.  It is not, and the
difference is the subject of this paper.

Consider the canonical failure.  A fitter parametrizes implied volatility
$\sigma(k)$ by some pleasant smooth curve, fits seven liquid strikes, and
reports residuals of a fraction of a basis point.  A risk system then
differentiates the fitted call prices twice with respect to strike ---
because that second derivative \emph{is} the risk-neutral probability
density, as \cref{sec:constraints} recalls --- and finds it negative between
two quoted strikes.  The fitted surface now prices a butterfly spread --- a
static portfolio of three vanilla options --- at a negative premium: the
fitter is offering free money, and reporting a negative probability.
Nothing failed numerically.  The parametrization simply did not know that a
smile is not free to be an arbitrary curve: it must be a probability
distribution seen through the Black formula, and the set of such curves has
a boundary that a generic parametric family crosses without noticing.

The standard remedy is to keep the pleasant curve and police the boundary:
evaluate a density proxy on a grid of strikes, penalize negativity, and hope
the optimizer terminates inside the feasible set.  This works, mostly, and
costs steadily.  The constraint is global and nonlinear in the parameters,
so it is re-checked at every strike and every iterate; the optimizer spends
its effort crawling along the boundary, where penalty gradients fight data
gradients; and in sparse or stressed markets --- precisely when a fitter is
most needed --- the boundary is where the solution wants to live.

This paper starts from a different question: \emph{in which representation
is validity free?}  That is, for which coordinates on the space of smiles is
every coordinate vector automatically arbitrage-free, so that nothing is
policed and the optimizer roams an open set?  The answer developed here is
the \emph{log-quantile-density} (LQD) model: write the log-forward return as
a monotone transport of a fixed logistic random variable, and parametrize
the \emph{logarithm of the transport speed} by a polynomial in the rank.
Positivity of the density is then structural --- an exponential is positive
--- the unit-mean condition is absorbed by one scalar shift, and the single
remaining condition, one inequality on one number, is removed by a final
change of variable.  The precise scope should be stated at once, because it
is the paper's first honesty obligation: the construction furnishes
\emph{unconstrained coordinates for a flexible exponential-tail class of
distributions}, not coordinates for the space of all smiles.  Atoms, bounded
support, density gaps, and Gaussian or faster tails lie outside the family;
\cref{sec:validity} says exactly what is covered and what is approximated.

\subsection{Three concerns, separated}

The design principle organizing the paper is that a smile fitter must manage
three separable concerns, and that the representation should separate them.

\begin{enumerate}[leftmargin=1.6em]
\item \textbf{Validity} --- no static arbitrage.  In the LQD representation
this is a property of the coordinates, not of the optimizer: every point of
the open parameter set is a genuine probability measure with unit mean
(\cref{sec:validity}).  There is no positivity penalty anywhere in the
construction, because there is nothing to penalize.
\item \textbf{Information} --- what the quoted strip actually determines.
The body of the distribution is pinned by quotes; the far tails are selected
by the basis, the regularization, and the endpoint coupling.  The paper says
this out loud and quantifies it: the moment strip and Lee slopes are exact
functions of two fitted numbers (\cref{sec:tails}), but those numbers are
priors dressed as parameters, and the honest deliverable at a traded strike
is the effective wing slope, which differs from the asymptotic limit by
economically large factors --- in either direction, on one and the same
fitted slice (\cref{fig:lee}).
\item \textbf{Convention} --- what the fitter chooses where the data is
silent.  Basis order, ridge weight, the optimization chart, the scope of
calendar enforcement: each is an explicit, auditable choice, orthogonal to
validity (\cref{sec:calibration,sec:calendar,sec:computation}).  A
convention defended as a theorem is a future incident; the paper labels its
conventions as such.
\end{enumerate}

One computational object carries all of this, and it plays the role in this
paper that a single unifying quantity plays in any theory worth the name.
Define the \emph{upper-share ledger}
\[
G(z)\;=\;\int_z^{\infty}e^{x(t)}\,\rho(t)\,\dd t ,
\]
the share of the (unit) forward carried by outcomes above the $z$-th
log-odds level; here $x$ is the transport map and $\rho$ the logistic
density, both defined in \cref{sec:model}.  Five computations that are
usually five pieces of software are five uses of this one array:
\begin{itemize}[leftmargin=1.4em]
\item \emph{pricing}: a call is the difference of two ledger entries
(\cref{sec:ledger});
\item \emph{differentiation}: the parameter Jacobian of every price follows
from an exact cancellation --- the strike root's motion drops out --- so one
quadrature pass differentiates the whole fit (\cref{sec:jacobian});
\item \emph{density}: the Breeden--Litzenberger density is one division of
arrays already on the grid (\cref{sec:ledger});
\item \emph{the variance swap}: the log-contract strike is the first moment
of the same build (\cref{sec:ledger});
\item \emph{calendar order}: two expiries are free of calendar arbitrage
precisely when their ledgers are pointwise ordered --- a theorem proved both
ways by conjugate duality (\cref{sec:calendar}).
\end{itemize}

\subsection{Contributions and route}

The ingredients are classical: Breeden and Litzenberger's density
identity~\citep{BreedenLitzenberger1978}, Lee's moment
formula~\citep{Lee2004}, quantile-function modeling in the statistical
tradition of Parzen~\citep{Parzen1979} and Gilchrist~\citep{Gilchrist2000},
the log-quantile-density transform studied by Petersen and
M\"uller~\citep{PetersenMuller2016}, and the log
contract~\citep{Neuberger1994}.  The contribution claimed is the
\emph{separation of concerns} --- validity structural, information audited,
convention explicit --- and the \emph{single-ledger computation} that makes
the separation cheap enough for production use.  Along the way the paper
proves the statements a referee would otherwise have to take on faith: the
finite-forward wall (\cref{prop:wall}), butterfly-freedom
(\cref{prop:noarb}), a universality proposition with its exclusions stated
(\cref{prop:universal}), closed-form tail-scale-to-Lee-slope maps
(\cref{prop:leeclosed}), exact at-the-money level, skew, and curvature with
the skew exhibited as a digital mismatch (\cref{prop:atm}), the
envelope-cancellation theorem behind the analytic Jacobian
(\cref{thm:envelope}), and the equivalence of calendar order with ledger
order (\cref{thm:ledgerorder}).

The paper is written for a reader who knows the Black--Scholes formula and
nothing else about smile modeling.  \Cref{sec:constraints} assembles the
constraint set; \cref{sec:heuristic} builds the representation from a
heuristic about ranks and returns; \cref{sec:model,sec:validity} define the
family and prove validity; \cref{sec:tails} treats tails and moments;
\cref{sec:ledger,sec:atm} develop the ledger and the at-the-money
microscope; \cref{sec:calibration,sec:jacobian,sec:calendar} cover
calibration, differentiation, and the calendar; \cref{sec:computation}
gives the numerical analysis, including the certificates by which the
discretization earns the continuous theorems; \cref{sec:examples} works
real-market and synthetic examples from one frozen data snapshot;
\cref{sec:related,sec:limits} position the work and state its limits.
Implementation detail is confined to \cref{app:implementation}; deferred
proofs to \cref{app:proofs}; data provenance to \cref{app:data}.
